\begin{abstract}
This paper asks a local-history question with national implications: how did James Wilkinson survive politically despite the scandals attached to his name, and what does his relationship with Thomas Jefferson have to do with Frederick, Maryland? Using a curated corpus of seventeen records, including fourteen primary documents and three Frederick-centered local sources, I argue that Wilkinson's relationship with Jefferson is best understood not as simple friendship but as a form of political intimacy built from usefulness, executive protection, and later retrospective esteem. The documentary pattern is striking. Only one coded record is primarily social, while ten fall into instrumental, protective, self-vindicating, or official-settlement modes. Frederick matters because it was one of the local places where that political survivability was processed into procedure. The paper's Frederick claim rests on a three-document primary spine --- the 1804 Frederick Town letter to Adams, the 1811 appeal from Frederick Town to Madison, and Madison's 1812 approval of the Frederick-town acquittal --- and is then thickened by a small set of later contextual sources. The computational package supports that interpretation rather than replacing it: Frederick-linked records make up thirty-five percent of the corpus, Frederick remains tied for the third-heaviest node even in the primary-only network, and the coded language remains weighted toward use, defense, and adjudication rather than simple warmth. The result is a more precise local claim: Wilkinson survived not because Jefferson cherished him in private, but because Jefferson's reassurance and Madison's ratification met a place where scandal could be turned into a survivable record.
\end{abstract}
